\begin{PSbox}[unbreakable]{obj}{اهداف یادگیری}

پس از تکمیل این ماژول، دانشجویان قادر خواهند بود:

\begin{enumerate}
\def\labelenumi{\arabic{enumi}.}
\tightlist
\item
  \lr{PMF}، \lr{PDF} و \lr{CDF} مشترک را تعریف و محاسبه کنند
\item
  توزیع‌های حاشیه‌ای را از توزیع‌های مشترک استخراج کنند
\item
  توزیع‌های شرطی را محاسبه کنند
\item
  استقلال دو متغیر تصادفی را آزمون کنند
\item
  امیدهای ریاضی شرطی را محاسبه کنند
\item
  توزیع‌های مشترک را در \lr{MATLAB} مصورسازی کنند
\end{enumerate}

\end{PSbox}

\PSsection{6.1}{مقدمه: چرا توزیع‌های مشترک را مطالعه کنیم؟}

سیستم‌های مهندسی واقعی شامل \textbf{چند کمیت نامعین برهم‌کنش‌کننده} هستند:

\begin{itemize}
\tightlist
\item
  ولتاژ و جریان در یک مدار نویزی
\item
  سیگنال و نویز در گیرنده
\item
  دما و مقاومت یک حسگر
\item
  کانال \lr{I} و کانال \lr{Q} یک سیگنال مخابراتی
\end{itemize}

درک \textbf{رفتار مشترک} چیزهایی را به ما می‌گوید که توزیع‌های حاشیه‌ای نمی‌توانند \lr{—} مانند همبستگی، وابستگی و رفتار شرطی.

\PSsection{6.2}{\lr{PMF} مشترک (حالت گسسته)}

\begin{PSbox}[unbreakable]{def}{تعریف}

برای متغیرهای تصادفی گسسته \lr{X}، \lr{Y}:\PSbr{}\[p_{X,Y}(x,y) = P(X = x, Y = y)\]

\end{PSbox}

\PSsubsection{ویژگی‌ها}

\begin{enumerate}
\def\labelenumi{\arabic{enumi}.}
\tightlist
\item
  \lr{p\textsubscript{X,Y}(x,y) \ensuremath{\ge} 0}
\item
  \lr{\ensuremath{\Sigma}\textsubscript{x} \ensuremath{\Sigma}\textsubscript{y} p\textsubscript{X,Y}(x,y) = 1}
\end{enumerate}

\PSsubsection{مصورسازی: جدول احتمال}

\textbf{مثال:} \lr{X =} تعداد ارسال‌های مجدد \PSnum{(0,1,2)}، \lr{Y =} دسته تأخیر تأیید (\lr{1=}سریع، \lr{2=}کند)

{\def\LTcaptype{none} % do not increment counter
\begin{longtable}[c]{@{}cccc@{}}
\toprule\noalign{}
\rowcolor{PSnavy}& \PSth{\lr{Y=1}} & \PSth{\lr{Y=2}} & \PSth{\lr{p\textsubscript{X}(x)}} \\
\midrule\noalign{}
\endhead
\bottomrule\noalign{}
\endlastfoot
\lr{X=0} & \PSnum{0.30} & \PSnum{0.10} & \PSnum{0.40} \\
\lr{X=1} & \PSnum{0.15} & \PSnum{0.20} & \PSnum{0.35} \\
\lr{X=2} & \PSnum{0.05} & \PSnum{0.20} & \PSnum{0.25} \\
\lr{p\textsubscript{Y}(y)} & \PSnum{0.50} & \PSnum{0.50} & \PSnum{1.00} \\
\end{longtable}
}

\PSsection{6.3}{\lr{PDF} مشترک (حالت پیوسته)}

\begin{PSbox}[unbreakable]{def}{تعریف}

برای متغیرهای تصادفی پیوسته \lr{X}، \lr{Y}، \lr{PDF} مشترک \lr{f\textsubscript{X,Y}(x,y)} در رابطه زیر صدق می‌کند:\PSbr{}\[P((X,Y) \in A) = \iint_A f_{X,Y}(x,y) \, dx \, dy\]

\end{PSbox}

\PSsubsection{ویژگی‌ها}

\begin{enumerate}
\def\labelenumi{\arabic{enumi}.}
\tightlist
\item
  \lr{f\textsubscript{X,Y}(x,y) \ensuremath{\ge} 0}
\item
  \lr{\ensuremath{\int}\ensuremath{\int} f\textsubscript{X,Y}(x,y) dx dy = 1}
\end{enumerate}

\PSsubsection{مصورسازی}

\lr{PDF} مشترک یک \textbf{رویه} روی صفحه \lr{(x,y)} است. احتمال \lr{=} حجم زیر رویه روی یک ناحیه.

\begin{PSbox}[unbreakable]{ex}{مثال مهندسی: گاوسی دوبعدی}

دو ولتاژ نویز همبسته \lr{(X, Y)} با همبستگی \lr{\ensuremath{\rho}}:

\[f_{X,Y}(x,y) = \frac{1}{2\pi\sigma_X\sigma_Y\sqrt{1-\rho^2}} \exp\left(-\frac{1}{2(1-\rho^2)}\left[\frac{x^2}{\sigma_X^2} - \frac{2\rho xy}{\sigma_X\sigma_Y} + \frac{y^2}{\sigma_Y^2}\right]\right)\]

\end{PSbox}

\PSsection{6.4}{\lr{CDF} مشترک}

\begin{PSbox}[unbreakable]{def}{تعریف}

\[F_{X,Y}(x,y) = P(X \leq x, Y \leq y)\]

\end{PSbox}

\PSsubsection{ویژگی‌ها}

\begin{enumerate}
\def\labelenumi{\arabic{enumi}.}
\tightlist
\item
  \lr{F\textsubscript{X,Y}(-\ensuremath{\infty}, y) = 0}، \lr{F\textsubscript{X,Y}(x, -\ensuremath{\infty}) = 0}
\item
  \lr{F\textsubscript{X,Y}(\ensuremath{\infty}, \ensuremath{\infty}) = 1}
\item
  نسبت به هر دو آرگومان نامنزولی است
\item
  \lr{f\textsubscript{X,Y}(x,y) = \ensuremath{\partial}\textsuperscript{2}F\textsubscript{X,Y}/(\ensuremath{\partial}x \ensuremath{\partial}y)}
\end{enumerate}

\PSsection{6.5}{توزیع‌های حاشیه‌ای}

\PSsubsection{به‌دست‌آوردن حاشیه‌ای‌ها از توزیع مشترک}

\textbf{گسسته:}\PSbr{}\[p_X(x) = \sum_y p_{X,Y}(x,y), \quad p_Y(y) = \sum_x p_{X,Y}(x,y)\]

\textbf{پیوسته:}\PSbr{}\[f_X(x) = \int_{-\infty}^{\infty} f_{X,Y}(x,y) \, dy, \quad f_Y(y) = \int_{-\infty}^{\infty} f_{X,Y}(x,y) \, dx\]

\PSsubsection{تفسیر بصری}

\lr{PDF} حاشیه‌ای \lr{f\textsubscript{X}(x)} همان \textbf{تصویر} (انتگرال) رویه مشترک روی محور \lr{x} است. به‌همین‌ترتیب \lr{f\textsubscript{Y}(y)} روی محور \lr{y} تصویر می‌شود.

\begin{PSbox}[unbreakable]{ex}{مثال مهندسی}

\lr{PDF} مشترک: \lr{f\textsubscript{X,Y}(x,y) = 2e\textsuperscript{\lr{(-x)e}}(-2y)}، \lr{x \ensuremath{\ge} 0}، \lr{y \ensuremath{\ge} 0}

حاشیه‌ای‌ها:

\begin{itemize}
\tightlist
\item
  \lr{f\textsubscript{X}(x) = \ensuremath{\int}\textsubscript{0}\textsuperscript{\ensuremath{\infty}} 2e\textsuperscript{\lr{(-x)e}}(-2y) dy = 2e\textsuperscript{-x} \ensuremath{\cdot} [1/2] = e\textsuperscript{-x} \ensuremath{\to} X \textasciitilde{} Exp(1)}
\item
  \lr{f\textsubscript{Y}(y) = \ensuremath{\int}\textsubscript{0}\textsuperscript{\ensuremath{\infty}} 2e\textsuperscript{\lr{(-x)e}}(-2y) dx = 2e\textsuperscript{-2y} \ensuremath{\cdot} [1] = 2e\textsuperscript{-2y} \ensuremath{\to} Y \textasciitilde{} Exp(2)}
\end{itemize}

\end{PSbox}

\PSsection{6.6}{توزیع‌های شرطی}

\PSsubsection{گسسته}

\[p_{X|Y}(x|y) = \frac{p_{X,Y}(x,y)}{p_Y(y)}\]

\PSsubsection{پیوسته}

\[f_{X|Y}(x|y) = \frac{f_{X,Y}(x,y)}{f_Y(y)}\]

\PSsubsection{تفسیر}

توزیع شرطی، \lr{X} را هنگامی توصیف می‌کند که \lr{Y = y} را \textbf{می‌دانیم}. این یک \lr{«}برش\lr{»} از توزیع مشترک در مقدار ثابت \lr{y} است که به‌گونه‌ای نرمال شده که انتگرال آن برابر 1 شود.

\begin{PSbox}[unbreakable]{ex}{مثال مهندسی: سیگنال به‌شرط وضعیت کانال}

توان دریافتی \lr{Y} به‌شرط وضعیت کانال \lr{X}:

\begin{itemize}
\tightlist
\item
  اگر کانال خوب باشد \lr{(X=1): Y \textasciitilde{} N(10, 1)}
\item
  اگر کانال ضعیف باشد \lr{(X=0): Y \textasciitilde{} N(2, 4)}
\item
  \lr{P(X=1) = 0.7}، \lr{P(X=0) = 0.3}
\end{itemize}

\lr{PDF}‌های شرطی رفتار گیرنده در هر وضعیت را توصیف می‌کنند.

\end{PSbox}

\PSsection{6.7}{استقلال دو متغیر تصادفی}

\begin{PSbox}[unbreakable]{def}{تعریف}

\lr{X} و \lr{Y} مستقل هستند اگر و تنها اگر:

\[f_{X,Y}(x,y) = f_X(x) \cdot f_Y(y) \quad \PSmtext{for all } x, y\]

(یا به‌طور معادل برای \lr{PMF}‌ها: \lr{p\textsubscript{X,Y}(x,y) = p\textsubscript{X}(x)\ensuremath{\cdot}p\textsubscript{Y}(y)})

\end{PSbox}

\PSsubsection{پیامدهای استقلال}

\begin{itemize}
\tightlist
\item
  \lr{f\textsubscript{X\textbar{}Y}(x\textbar{}y) = f\textsubscript{X}(x)} (دانستن \lr{Y} مقدار \lr{X} را تغییر نمی‌دهد)
\item
  \lr{E[XY] = E[X]E[Y]}
\item
  \lr{Var(X + Y) = Var(X) + Var(Y)}
\end{itemize}

\PSsubsection{آزمون استقلال}

بررسی کنید: آیا توزیع مشترک به حاصل‌ضرب یک تابع فقط از \lr{x} و یک تابع فقط از \lr{y} تجزیه می‌شود؟

\begin{itemize}
\tightlist
\item
  \lr{f\textsubscript{X,Y}(x,y) = 2e\textsuperscript{\lr{(-x)\ensuremath{\cdot}e}}(-2y) = [e\textsuperscript{-x}]\ensuremath{\cdot}[2e\textsuperscript{-2y}] \PSyes{}} مستقل!
\item
  \lr{f\textsubscript{X,Y}(x,y) = x + y} برای \lr{0\ensuremath{\le}x,y\ensuremath{\le}1 \ensuremath{\to}} تجزیه نمی‌شود \lr{\ensuremath{\to}} مستقل نیست
\end{itemize}

\begin{PSbox}[unbreakable]{ex}{مثال مهندسی: سیگنال و نویز مستقل}

سیگنال ارسالی \lr{S} و نویز کانال \lr{N} بر اساس فرض فیزیکی مستقل هستند:\PSbr{}\lr{f\textsubscript{S,N}(s,n) = f\textsubscript{S}(s)\ensuremath{\cdot}f\textsubscript{N}(n)}

این فرض استقلال برای طراحی سیستم مخابراتی بنیادی است.

\end{PSbox}

\PSsection{6.8}{امید ریاضی شرطی}

\begin{PSbox}[unbreakable]{def}{تعریف}

\[E[X|Y=y] = \begin{cases} \sum_x x \cdot p_{X|Y}(x|y) & \PSmtext{discrete} \\ \int_{-\infty}^{\infty} x \cdot f_{X|Y}(x|y) \, dx & \PSmtext{continuous} \end{cases}\]

\end{PSbox}

\PSsubsection{تفسیر}

\lr{E[X\textbar{}Y=y]} \textbf{بهترین پیش‌بینی} \lr{X} است به‌شرط این‌که \lr{Y = y} را مشاهده کرده‌ایم (به معنای میانگین مربعات).

\begin{PSbox}[unbreakable]{ex}{مثال مهندسی: مقاومت وابسته به دما}

مقدار یک مقاومت \lr{R} به دما \lr{T} وابسته است. اگر \lr{R\textbar{}T \textasciitilde{} N(1000 + 0.5T, 4)}:

\lr{E[R\textbar{}T=50\ensuremath{^{\circ}}C] = 1000 + 0.5(50) = 1025\ensuremath{\Omega}}

با دانستن دما، می‌توانیم مقاومت را پیش‌بینی کنیم.

\end{PSbox}

\PSsection{6.9}{مثال‌های \lr{MATLAB}}

\begin{PSbox}[unbreakable]{ex}{مثال 1: مصورسازی \lr{PMF} مشترک}

\tcbset{PScodemode/.style={unbreakable}}

\begin{Shaded}
\begin{Highlighting}[]
\CommentTok{\%\% Joint PMF: Retransmissions vs Delay Category}
\VariableTok{pXY} \OperatorTok{=}\NormalTok{ [}\FloatTok{0.30} \FloatTok{0.10}\OperatorTok{;} \FloatTok{0.15} \FloatTok{0.20}\OperatorTok{;} \FloatTok{0.05} \FloatTok{0.20}\NormalTok{]}\OperatorTok{;}
\VariableTok{x} \OperatorTok{=} \FloatTok{0}\OperatorTok{:}\FloatTok{2}\OperatorTok{;} \VariableTok{y} \OperatorTok{=} \FloatTok{1}\OperatorTok{:}\FloatTok{2}\OperatorTok{;}

\VariableTok{figure}\OperatorTok{;}
\VariableTok{bar3}\NormalTok{(}\VariableTok{pXY}\NormalTok{)}\OperatorTok{;}
\VariableTok{xlabel}\NormalTok{(}\SpecialStringTok{\textquotesingle{}Delay Category Y\textquotesingle{}}\NormalTok{)}\OperatorTok{;} \VariableTok{ylabel}\NormalTok{(}\SpecialStringTok{\textquotesingle{}Retransmissions X\textquotesingle{}}\NormalTok{)}\OperatorTok{;}
\VariableTok{zlabel}\NormalTok{(}\SpecialStringTok{\textquotesingle{}P(X=x, Y=y)\textquotesingle{}}\NormalTok{)}\OperatorTok{;}
\VariableTok{title}\NormalTok{(}\SpecialStringTok{\textquotesingle{}Joint PMF\textquotesingle{}}\NormalTok{)}\OperatorTok{;}

\CommentTok{\% Marginals}
\VariableTok{pX} \OperatorTok{=} \VariableTok{sum}\NormalTok{(}\VariableTok{pXY}\OperatorTok{,} \FloatTok{2}\NormalTok{)}\OperatorTok{\textquotesingle{};}   \CommentTok{\% Sum over Y}
\VariableTok{pY} \OperatorTok{=} \VariableTok{sum}\NormalTok{(}\VariableTok{pXY}\OperatorTok{,} \FloatTok{1}\NormalTok{)}\OperatorTok{;}    \CommentTok{\% Sum over X}
\VariableTok{fprintf}\NormalTok{(}\SpecialStringTok{\textquotesingle{}Marginal of X: [\%.2f \%.2f \%.2f]\textbackslash{}n\textquotesingle{}}\OperatorTok{,} \VariableTok{pX}\NormalTok{)}\OperatorTok{;}
\VariableTok{fprintf}\NormalTok{(}\SpecialStringTok{\textquotesingle{}Marginal of Y: [\%.2f \%.2f]\textbackslash{}n\textquotesingle{}}\OperatorTok{,} \VariableTok{pY}\NormalTok{)}\OperatorTok{;}
\end{Highlighting}
\end{Shaded}

\end{PSbox}

\begin{PSbox}[unbreakable]{ex}{مثال 2: مصورسازی گاوسی دوبعدی}

\tcbset{PScodemode/.style={unbreakable}}

\begin{Shaded}
\begin{Highlighting}[]
\CommentTok{\%\% Joint Gaussian PDF with Correlation}
\VariableTok{mu} \OperatorTok{=}\NormalTok{ [}\FloatTok{0} \FloatTok{0}\NormalTok{]}\OperatorTok{;} \VariableTok{sigma\_x} \OperatorTok{=} \FloatTok{1}\OperatorTok{;} \VariableTok{sigma\_y} \OperatorTok{=} \FloatTok{1.5}\OperatorTok{;} \VariableTok{rho} \OperatorTok{=} \FloatTok{0.7}\OperatorTok{;}
\VariableTok{Sigma} \OperatorTok{=}\NormalTok{ [}\VariableTok{sigma\_x}\OperatorTok{\^{}}\FloatTok{2}\OperatorTok{,} \VariableTok{rho}\OperatorTok{*}\VariableTok{sigma\_x}\OperatorTok{*}\VariableTok{sigma\_y}\OperatorTok{;} \VariableTok{rho}\OperatorTok{*}\VariableTok{sigma\_x}\OperatorTok{*}\VariableTok{sigma\_y}\OperatorTok{,} \VariableTok{sigma\_y}\OperatorTok{\^{}}\FloatTok{2}\NormalTok{]}\OperatorTok{;}

\NormalTok{[}\VariableTok{X}\OperatorTok{,} \VariableTok{Y}\NormalTok{] }\OperatorTok{=} \VariableTok{meshgrid}\NormalTok{(}\VariableTok{linspace}\NormalTok{(}\OperatorTok{{-}}\FloatTok{4}\OperatorTok{,} \FloatTok{4}\OperatorTok{,} \FloatTok{100}\NormalTok{)}\OperatorTok{,} \VariableTok{linspace}\NormalTok{(}\OperatorTok{{-}}\FloatTok{5}\OperatorTok{,} \FloatTok{5}\OperatorTok{,} \FloatTok{100}\NormalTok{))}\OperatorTok{;}
\VariableTok{XY} \OperatorTok{=}\NormalTok{ [}\VariableTok{X}\NormalTok{(}\OperatorTok{:}\NormalTok{) }\VariableTok{Y}\NormalTok{(}\OperatorTok{:}\NormalTok{)]}\OperatorTok{;}
\VariableTok{Z} \OperatorTok{=} \VariableTok{mvnpdf}\NormalTok{(}\VariableTok{XY}\OperatorTok{,} \VariableTok{mu}\OperatorTok{,} \VariableTok{Sigma}\NormalTok{)}\OperatorTok{;}
\VariableTok{Z} \OperatorTok{=} \VariableTok{reshape}\NormalTok{(}\VariableTok{Z}\OperatorTok{,} \VariableTok{size}\NormalTok{(}\VariableTok{X}\NormalTok{))}\OperatorTok{;}

\VariableTok{figure}\OperatorTok{;}
\VariableTok{subplot}\NormalTok{(}\FloatTok{1}\OperatorTok{,}\FloatTok{2}\OperatorTok{,}\FloatTok{1}\NormalTok{)}\OperatorTok{;}
\VariableTok{surf}\NormalTok{(}\VariableTok{X}\OperatorTok{,} \VariableTok{Y}\OperatorTok{,} \VariableTok{Z}\OperatorTok{,} \SpecialStringTok{\textquotesingle{}EdgeColor\textquotesingle{}}\OperatorTok{,} \SpecialStringTok{\textquotesingle{}none\textquotesingle{}}\NormalTok{)}\OperatorTok{;} \VariableTok{view}\NormalTok{(}\FloatTok{30}\OperatorTok{,} \FloatTok{40}\NormalTok{)}\OperatorTok{;}
\VariableTok{xlabel}\NormalTok{(}\SpecialStringTok{\textquotesingle{}X\textquotesingle{}}\NormalTok{)}\OperatorTok{;} \VariableTok{ylabel}\NormalTok{(}\SpecialStringTok{\textquotesingle{}Y\textquotesingle{}}\NormalTok{)}\OperatorTok{;} \VariableTok{zlabel}\NormalTok{(}\SpecialStringTok{\textquotesingle{}f\_\{X,Y\}(x,y)\textquotesingle{}}\NormalTok{)}\OperatorTok{;}
\VariableTok{title}\NormalTok{(}\VariableTok{sprintf}\NormalTok{(}\SpecialStringTok{\textquotesingle{}Joint Gaussian (ρ = \%.1f)\textquotesingle{}}\OperatorTok{,} \VariableTok{rho}\NormalTok{))}\OperatorTok{;}

\VariableTok{subplot}\NormalTok{(}\FloatTok{1}\OperatorTok{,}\FloatTok{2}\OperatorTok{,}\FloatTok{2}\NormalTok{)}\OperatorTok{;}
\VariableTok{contour}\NormalTok{(}\VariableTok{X}\OperatorTok{,} \VariableTok{Y}\OperatorTok{,} \VariableTok{Z}\OperatorTok{,} \FloatTok{20}\NormalTok{)}\OperatorTok{;}
\VariableTok{xlabel}\NormalTok{(}\SpecialStringTok{\textquotesingle{}X\textquotesingle{}}\NormalTok{)}\OperatorTok{;} \VariableTok{ylabel}\NormalTok{(}\SpecialStringTok{\textquotesingle{}Y\textquotesingle{}}\NormalTok{)}\OperatorTok{;}
\VariableTok{title}\NormalTok{(}\SpecialStringTok{\textquotesingle{}Contour Plot\textquotesingle{}}\NormalTok{)}\OperatorTok{;}
\VariableTok{axis} \VariableTok{equal}\OperatorTok{;}
\end{Highlighting}
\end{Shaded}

\end{PSbox}

\begin{PSbox}[unbreakable]{ex}{مثال 3: آزمون استقلال}

\tcbset{PScodemode/.style={unbreakable}}

\begin{Shaded}
\begin{Highlighting}[]
\CommentTok{\%\% Independence Check via Simulation}
\VariableTok{N} \OperatorTok{=} \FloatTok{100000}\OperatorTok{;}

\CommentTok{\% Case 1: Independent (signal and noise)}
\VariableTok{S} \OperatorTok{=} \VariableTok{randn}\NormalTok{(}\FloatTok{1}\OperatorTok{,} \VariableTok{N}\NormalTok{)}\OperatorTok{;}               \CommentTok{\% Signal}
\VariableTok{Noise} \OperatorTok{=} \FloatTok{0.5}\OperatorTok{*}\VariableTok{randn}\NormalTok{(}\FloatTok{1}\OperatorTok{,} \VariableTok{N}\NormalTok{)}\OperatorTok{;}       \CommentTok{\% Independent noise}
\VariableTok{R} \OperatorTok{=} \VariableTok{S} \OperatorTok{+} \VariableTok{Noise}\OperatorTok{;}                  \CommentTok{\% Received}

\CommentTok{\% Case 2: Dependent (voltage and current through same resistor)}
\VariableTok{V} \OperatorTok{=} \VariableTok{randn}\NormalTok{(}\FloatTok{1}\OperatorTok{,} \VariableTok{N}\NormalTok{)}\OperatorTok{;}               \CommentTok{\% Voltage}
\VariableTok{I} \OperatorTok{=} \VariableTok{V} \OperatorTok{/} \FloatTok{50} \OperatorTok{+} \FloatTok{0.01}\OperatorTok{*}\VariableTok{randn}\NormalTok{(}\FloatTok{1}\OperatorTok{,}\VariableTok{N}\NormalTok{)}\OperatorTok{;} \CommentTok{\% Current (dependent on V!)}

\VariableTok{fprintf}\NormalTok{(}\SpecialStringTok{\textquotesingle{}Independent (S, Noise): corr = \%.4f\textbackslash{}n\textquotesingle{}}\OperatorTok{,} \VariableTok{corr}\NormalTok{(}\VariableTok{S}\OperatorTok{\textquotesingle{},} \VariableTok{Noise}\OperatorTok{\textquotesingle{}}\NormalTok{))}\OperatorTok{;}
\VariableTok{fprintf}\NormalTok{(}\SpecialStringTok{\textquotesingle{}Dependent (V, I): corr = \%.4f\textbackslash{}n\textquotesingle{}}\OperatorTok{,} \VariableTok{corr}\NormalTok{(}\VariableTok{V}\OperatorTok{\textquotesingle{},} \VariableTok{I}\NormalTok{))}\OperatorTok{;}
\end{Highlighting}
\end{Shaded}

\end{PSbox}

\begin{PSbox}[unbreakable]{ex}{مثال 4: توزیع شرطی}

\tcbset{PScodemode/.style={unbreakable}}

\begin{Shaded}
\begin{Highlighting}[]
\CommentTok{\%\% Conditional PDF: Slice of Bivariate Gaussian}
\VariableTok{rho} \OperatorTok{=} \FloatTok{0.8}\OperatorTok{;} \VariableTok{N} \OperatorTok{=} \FloatTok{100000}\OperatorTok{;}
\VariableTok{X} \OperatorTok{=} \VariableTok{randn}\NormalTok{(}\FloatTok{1}\OperatorTok{,} \VariableTok{N}\NormalTok{)}\OperatorTok{;}
\VariableTok{Y} \OperatorTok{=} \VariableTok{rho}\OperatorTok{*}\VariableTok{X} \OperatorTok{+} \VariableTok{sqrt}\NormalTok{(}\FloatTok{1}\OperatorTok{{-}}\VariableTok{rho}\OperatorTok{\^{}}\FloatTok{2}\NormalTok{)}\OperatorTok{*}\VariableTok{randn}\NormalTok{(}\FloatTok{1}\OperatorTok{,} \VariableTok{N}\NormalTok{)}\OperatorTok{;}  \CommentTok{\% Y|X \textasciitilde{} N(rho*x, 1{-}rho\^{}2)}

\CommentTok{\% Conditional: Y given X ≈ 1 (take samples where 0.9 \textless{} X \textless{} 1.1)}
\VariableTok{mask} \OperatorTok{=}\NormalTok{ (}\VariableTok{X} \OperatorTok{\textgreater{}} \FloatTok{0.9}\NormalTok{) }\OperatorTok{\&}\NormalTok{ (}\VariableTok{X} \OperatorTok{\textless{}} \FloatTok{1.1}\NormalTok{)}\OperatorTok{;}
\VariableTok{Y\_cond} \OperatorTok{=} \VariableTok{Y}\NormalTok{(}\VariableTok{mask}\NormalTok{)}\OperatorTok{;}

\VariableTok{figure}\OperatorTok{;}
\VariableTok{histogram}\NormalTok{(}\VariableTok{Y\_cond}\OperatorTok{,} \FloatTok{50}\OperatorTok{,} \SpecialStringTok{\textquotesingle{}Normalization\textquotesingle{}}\OperatorTok{,} \SpecialStringTok{\textquotesingle{}pdf\textquotesingle{}}\NormalTok{)}\OperatorTok{;}
\VariableTok{hold} \VariableTok{on}\OperatorTok{;}
\VariableTok{y} \OperatorTok{=} \VariableTok{linspace}\NormalTok{(}\OperatorTok{{-}}\FloatTok{3}\OperatorTok{,} \FloatTok{4}\OperatorTok{,} \FloatTok{200}\NormalTok{)}\OperatorTok{;}
\VariableTok{plot}\NormalTok{(}\VariableTok{y}\OperatorTok{,} \VariableTok{normpdf}\NormalTok{(}\VariableTok{y}\OperatorTok{,} \VariableTok{rho}\OperatorTok{*}\FloatTok{1}\OperatorTok{,} \VariableTok{sqrt}\NormalTok{(}\FloatTok{1}\OperatorTok{{-}}\VariableTok{rho}\OperatorTok{\^{}}\FloatTok{2}\NormalTok{))}\OperatorTok{,} \SpecialStringTok{\textquotesingle{}r\textquotesingle{}}\OperatorTok{,} \SpecialStringTok{\textquotesingle{}LineWidth\textquotesingle{}}\OperatorTok{,} \FloatTok{2}\NormalTok{)}\OperatorTok{;}
\VariableTok{xlabel}\NormalTok{(}\SpecialStringTok{\textquotesingle{}Y\textquotesingle{}}\NormalTok{)}\OperatorTok{;} \VariableTok{ylabel}\NormalTok{(}\SpecialStringTok{\textquotesingle{}PDF\textquotesingle{}}\NormalTok{)}\OperatorTok{;}
\VariableTok{title}\NormalTok{(}\SpecialStringTok{\textquotesingle{}f\_\{Y|X\}(y|X≈1) for Bivariate Gaussian\textquotesingle{}}\NormalTok{)}\OperatorTok{;}
\VariableTok{legend}\NormalTok{(}\SpecialStringTok{\textquotesingle{}Empirical\textquotesingle{}}\OperatorTok{,} \VariableTok{sprintf}\NormalTok{(}\SpecialStringTok{\textquotesingle{}N(\%.1f, \%.2f)\textquotesingle{}}\OperatorTok{,} \VariableTok{rho}\OperatorTok{,} \FloatTok{1}\OperatorTok{{-}}\VariableTok{rho}\OperatorTok{\^{}}\FloatTok{2}\NormalTok{))}\OperatorTok{;}
\end{Highlighting}
\end{Shaded}

\end{PSbox}

\PSsection{6.10}{مسایل تمرینی}

\begin{PSbox}[unbreakable]{prob}{مسیله 1}

\lr{PDF} مشترک: \lr{f\textsubscript{X,Y}(x,y) = 6(1-y)} برای \lr{0 \ensuremath{\le} x \ensuremath{\le} y \ensuremath{\le} 1} و در جای دیگر صفر.

\begin{enumerate}
\def\labelenumi{(\alph{enumi})}
\tightlist
\item
  نرمال‌سازی را بررسی کنید. \lr{(b)} حاشیه‌ای‌ها را بیابید. \lr{(c)} آیا \lr{X} و \lr{Y} مستقل هستند؟ \lr{(d) P(X \textless{} 0.5, Y \textless{} 0.5)} را بیابید.
\end{enumerate}

\PSsolution{} 

\begin{enumerate}
\def\labelenumi{(\alph{enumi})}
\tightlist
\item
  \lr{\ensuremath{\int}\textsubscript{0}\textsuperscript{1} \ensuremath{\int}\textsubscript{0}\textsuperscript{y} 6(1-y) dx dy = \ensuremath{\int}\textsubscript{0}\textsuperscript{1} 6y(1-y) dy = 6[y\textsuperscript{2}/2 - y\textsuperscript{3}/3]\textsubscript{0}\textsuperscript{1} = 6(1/2 - 1/3) = 1 \PSyes{}}
\item
  \lr{f\textsubscript{X}(x) = \ensuremath{\int}\textsubscript{x}\textsuperscript{1} 6(1-y) dy = 6[(1-y\textsuperscript{2}/2) - (y)]\textsubscript{x}\textsuperscript{1} = 3(1-x)\textsuperscript{2}}؛ \lr{f\textsubscript{Y}(y) = 6y(1-y)}
\item
  \lr{f\textsubscript{X,Y} \ensuremath{\neq} f\textsubscript{X} \ensuremath{\cdot} f\textsubscript{Y} \ensuremath{\to}} مستقل نیستند (همچنین تکیه‌گاه مثلثی است، نه مستطیلی)
\item
  \lr{P(X\textless{}0.5, Y\textless{}0.5) = \ensuremath{\int}\textsubscript{0}\textsuperscript{0}.5 \ensuremath{\int}\textsubscript{0}\textsuperscript{y} 6(1-y) dx dy = \ensuremath{\int}\textsubscript{0}\textsuperscript{0}.5 6y(1-y) dy = 6[y\textsuperscript{2}/2 - y\textsuperscript{3}/3]\textsubscript{0}\textsuperscript{0}.5 = 0.500}
\end{enumerate}

\end{PSbox}

\begin{PSbox}[unbreakable]{prob}{مسیله 2}

دو حسگر یک سیگنال را با نویز مستقل اندازه‌گیری می‌کنند. \lr{X = S + N\textsubscript{1}}، \lr{Y = S + N\textsubscript{2}} که در آن \lr{S=5V}، \lr{N\textsubscript{1}\textasciitilde{}N(0,1)}، \lr{N\textsubscript{2}\textasciitilde{}N(0,4)}. آیا \lr{X} و \lr{Y} مستقل هستند؟ \lr{E[X]}، \lr{E[Y]} را بیابید و توضیح دهید چرا \lr{Cov(X,Y) \ensuremath{\neq} 0}.

\PSsolution{} \lr{X} و \lr{Y} مستقل نیستند زیرا در \lr{S} مشترک‌اند (هرچند نویز مستقل است). \lr{E[X] = E[Y] = 5. Cov(X,Y) = Cov(S+N\textsubscript{1}, S+N\textsubscript{2}) = Var(S) + Cov(S,N\textsubscript{1}) + Cov(N\textsubscript{2},S) + Cov(N\textsubscript{1},N\textsubscript{2})}. چون \lr{S} ثابت است، \lr{Var(S)=0}، بنابراین اگر \lr{S} تصادفی بود، آن‌ها همبسته می‌شدند. با \lr{S} ثابت، \lr{X} و \lr{Y} در این حالت در واقع مستقل هستند. اگر \lr{S} تصادفی با \lr{Var(S) = \ensuremath{\sigma}\textsubscript{S}\textsuperscript{2}} باشد، آنگاه \lr{Cov(X,Y) = \ensuremath{\sigma}\textsubscript{S}\textsuperscript{2}}.

\end{PSbox}

\begin{PSbox}[unbreakable]{prob}{مسیله 3}

با توجه به \lr{PMF} مشترک در جدول بخش \PSnum{6.2}، موارد زیر را بیابید: \lr{(a) P(X\ensuremath{\le}1, Y=2)}، \lr{(b) f\textsubscript{X\textbar{}Y}(x\textbar{}Y=1)}، \lr{(c) E[X\textbar{}Y=1]}.

\PSsolution{} 

\begin{enumerate}
\def\labelenumi{(\alph{enumi})}
\tightlist
\item
  \lr{P(X\ensuremath{\le}1, Y=2) = p(0,2) + p(1,2) = 0.10 + 0.20 = 0.30}
\item
  \lr{p\textsubscript{X\textbar{}Y}(x\textbar{}1) = p(x,1)/p\textsubscript{Y}(1): p(0\textbar{}1)=0.30/0.50=0.60}، \lr{p(1\textbar{}1)=0.15/0.50=0.30}، \lr{p(2\textbar{}1)=0.05/0.50=0.10}
\item
  \lr{E[X\textbar{}Y=1] = 0(0.60) + 1(0.30) + 2(0.10) = 0.50}
\end{enumerate}

\emph{ماژول بعدی: {ماژول \lr{7 —} توابع دو متغیر تصادفی}}\PSbr{}

\end{PSbox}
